\section{ภาคผนวก}
\subsection{สคริปต์ MATLAB สำหรับการวิเคราะห์หาค่าพารามิเตอร์ของระบบ}

% สลับการตั้งค่าภาษาชั่วคราวให้รองรับ Syntax ของ MATLAB
\lstset{language=Matlab}

\begin{lstlisting}[caption={MATLAB Script for Spring Constant (K) and Damping Coefficient (C) Data Preparation}]
clear; clc; close all;

%% 1. Spring Constant (K) Calculation via Full OLS
% Initialize arrays for measured Force (N) and Displacement (m)
force_data = [0.000, 0.372, 0.743, 1.116, 1.488, 1.860];
disp_data  = [0.00000, 0.01222, 0.02429, 0.03735, 0.05007, 0.06407];

% Perform 1st-degree polynomial curve fitting (Linear Regression)
% The function returns coefficients: p(1) = slope (K), p(2) = intercept (c)
p_coeffs = polyfit(disp_data, force_data, 1);
K_estimate = p_coeffs(1);
c_offset   = p_coeffs(2);

fprintf('=== System Parameter Estimation ===\n');
fprintf('Estimated Spring Constant (K): %.2f N/m\n', K_estimate);
fprintf('Systematic Measurement Offset (c): %.4f N\n', c_offset);

%% 2. Damping Coefficient (C) Identification Setup
% Load free vibration trial data from the pre-processed .mat file
% The variables must match the workspace expected by Simulink Design Optimization
load('ScenarioB_FreeVibration.mat', 't_ms', 'filtered_disp');

% Convert timestamp to seconds for proper continuous domain scaling
time_sec = t_ms / 1000.0;

fprintf('\nData successfully loaded for Simulink Parameter Estimator.\n');
% Note: The non-linear least squares estimation for C is executed 
% interactively via the Simulink Parameter Estimator application.
\end{lstlisting}